
\chapter{Rule Discovery Results (Goal 1)}\label{ch:goal_1_rule_extraction}

This chapter reports the output of the three rule-learning algorithms applied to the full historical datasets. The primary evaluation question for Goal 1 is: does CarenR's distributional subgroup discovery extract statistically validated, non-trivial patterns from the data? Subsidiary questions are: do RIPPER and M5Rules converge on the same variables through different optimisation criteria? And are the discovered patterns genuine within-era signal or era-membership artefacts? Section~\ref{sec:carenr_distributional_rules} presents the CarenR rules; Section~\ref{sec:loess_validation_distinguishing_signal} validates them using LOESS detrending; Sections~\ref{sec:ripper_classification_rules}--\ref{sec:m5rules_piecewise_linear_models} present RIPPER and M5Rules; Sections~\ref{sec:crossalgorithm_feature_agreement}--\ref{sec:summary_goal_1_assessment} provide cross-algorithm comparison, feature frequency, and rule set structure analysis.


\section{CarenR: Distributional Rules}\label{sec:carenr_distributional_rules}

After applying the significance threshold (KS p $\leq$ 0.10) and minimum support filter ($\geq$ 20\% of training observations), CarenR retained 48 rules for RDD and 20 rules for RVV. Tables~\ref{tab:ch4a_1} and~\ref{tab:ch4a_2} show the five most significant rules per region; the full top-10 tables are in Appendix~\ref{app:complete_rule_listings} (Tables~\ref{tab:appa_top10_rdd} and~\ref{tab:appa_top10_rvv}). For each rule, the table records the statistical evidence (p-value, support), the magnitude of the discovered distributional shift (mean deviation from trend), and the conditions that define the subgroup. One sentence of domain context is provided per rule to confirm that the direction of the discovered pattern is consistent with established knowledge, serving as a face-validity check on the method's output rather than as a primary finding.


\subsection{Douro (RDD): Top 5 Rules}\label{sub:douro_rdd_top_10}

\begin{table}[!ht]
  \centering
  \caption{Top 5 CarenR rules for the Douro (RDD), sorted by KS $p$-value. Full top-10 in Table~\ref{tab:appa_top10_rdd}.}\label{tab:ch4a_1}
  \small
  \begin{adjustbox}{max width=\textwidth}
    \begin{tabularx}{\textwidth}{p{0.4cm}|X|>{\centering\arraybackslash}p{0.9cm}|p{1.3cm}|p{1.2cm}|p{1.2cm}}
  \toprule
  \# & Conditions & Dir. & Median $\Delta$ \newline (mhl) & Support & p-value \\
  \midrule
  1 & \textbf{Wet days post-Flowering [0--3.3d] AND Wet days post-Maturity [0--3.5d]} \newline \textit{Dry post-flowering conditions reduce disease pressure and promote fruit set in the Douro.} & {\footnotesize ABOVE} & +121 & 28.1\% & 0.0043 \\
  2 & \textbf{LAI at Flowering [0.74--0.88 m$^2$/m$^2$] AND Soil water post-Flowering [390--404 mm]} \newline \textit{Elevated canopy vigour with high soil moisture at flowering drives vegetative-over-reproductive growth.} & {\footnotesize BELOW} & -131 & 20.2\% & 0.0084 \\
  3 & \textbf{Wet days post-Flowering [0--3.3d] AND Wet days post-Maturity [0--3.5d] AND Heat days post-Flowering [0--0.6d]} \newline \textit{Refines Rule 1: dry and thermally moderate post-flowering window.} & {\footnotesize ABOVE} & +113 & 25.8\% & 0.0089 \\
  4 & \textbf{Soil water pre-Harvest [74--101 mm] AND Wet days post-Harvest prev.yr [0--2.3d] AND Wet days post-Maturity [0--3.5d]} \newline \textit{Multi-year rule: moderate soil water stress with dry previous-year post-harvest window captures a carry-over vine reserve mechanism.} & {\footnotesize ABOVE} & +192 & 21.3\% & 0.0095 \\
  5 & \textbf{Wet days post-Flowering [0--3.3d] AND Heat days pre-Maturity [0--3.7d]} \newline \textit{Dry flowering with bounded pre-veraison heat — the ``warm but not extreme'' pattern.} & {\footnotesize ABOVE} & +115 & 24.7\% & 0.0110 \\
  \bottomrule
\end{tabularx}
  \end{adjustbox}

\end{table}


Key structural observations from the RDD rule set: (1) 31 of 48 rules point to ABOVE-trend production, an asymmetry that reflects the algorithm's finding that the positive signal in this dataset is structurally more consistent than the negative signal; the 17 below-trend rules require an unusual combination of high LAI with high soil moisture that occurs in approximately 20\% of years. (2) The rule set is organised around three recurring features, post-flowering wet days, pre-harvest soil water, and heat-stress day counts, with later rules being refinements of earlier ones through additional conditions. (3) The most significant rule (Rule 1) has support of 28.1\%, covering approximately one in four years in the historical record, confirming it describes a recurrent rather than exceptional pattern. Rule 4 produces the largest above-trend shift in the RDD set (+192 mhl, support 21.3\%), combining moderate pre-harvest soil water with dry post-harvest conditions across two years; Rule 2 is the largest below-trend rule (-131 mhl, support 20.2\%), requiring elevated canopy vigour with high soil moisture at flowering.


\subsection{Vinho Verde (RVV): Top 5 Rules}\label{sub:vinho_verde_rvv_top}

\begin{table}[!ht]
  \centering
  \caption{Top 5 CarenR rules for the Vinho Verde (RVV). Rule 4 is the only above-trend rule. Full top-10 in Table~\ref{tab:appa_top10_rvv}.}\label{tab:ch4a_2}
  \small
  \begin{adjustbox}{max width=\textwidth}
    \begin{tabularx}{\textwidth}{p{0.4cm}|X|>{\centering\arraybackslash}p{0.9cm}|p{1.3cm}|p{1.2cm}|p{1.2cm}}
  \toprule
  \# & Conditions & Dir. & Median $\Delta$ \newline (mhl) & Support & p-value \\
  \midrule
  1 & \textbf{Soil water pre-Harvest [61--97 mm] AND Heat days pre-Maturity [0--1.9d] AND Heat days post-Maturity [0--2.4d] AND Heat days post-Flowering [0--1.1d]} \newline \textit{Low soil water deficit with absent heat stress — in the Atlantic RVV climate this is a cool-moist pattern, not a beneficial stress pattern.} & {\footnotesize BELOW} & -338 & 21.2\% & 0.0198 \\
  2 & \textbf{Temp at Flowering [13.8--17.1$^\circ$C] AND Temp post-Flowering [12.6--17.5$^\circ$C] AND Wet-soil days post-Flowering prev.yr [17.5--20d]} \newline \textit{Cool flowering combined with previous-year soil saturation — a two-year carry-over pattern.} & {\footnotesize BELOW} & -268 & 21.2\% & 0.0216 \\
  3 & \textbf{Soil water post-Flowering [143--234 mm] AND Heat days post-Maturity [0--2.4d]} \newline \textit{Highest-support below-trend rule: moist post-flowering with low heat at maturity.} & {\footnotesize BELOW} & -365 & 27.5\% & 0.0294 \\
  4 & \textbf{Temp post-Budburst prev.yr [11.2--13.5$^\circ$C] AND Cold days at Flowering [0--3.4d] AND Heat days post-Flowering [0--1.1d]} \newline \textit{The only above-trend rule: warm-but-not-extreme previous-year budburst with mild current-year flowering. Largest absolute deviation in either region (+311 mhl).} & {\footnotesize ABOVE} & +311 & 25.0\% & 0.0298 \\
  5 & \textbf{Temp post-Flowering [12.6--17.5$^\circ$C] AND Heat days post-Maturity [0--2.4d] AND Wet-soil days post-Flowering prev.yr [17.5--20d]} \newline \textit{Cool post-flowering temperature with previous-year soil saturation carry-over.} & {\footnotesize BELOW} & -292 & 23.8\% & 0.0299 \\
  \bottomrule
\end{tabularx}
  \end{adjustbox}

\end{table}


Key structural observations from the RVV rule set: (1) 16 of 20 rules point to BELOW-trend production, the inverse asymmetry from RDD. In a declining series, the majority of historically exceptional subgroups are years that fall below the already-declining trend. (2) Rule 4 (one of four above-trend rules, +311 mhl, support 25\%) represents the largest absolute distributional shift in either region; CarenR has identified a well-defined subgroup corresponding to the best RVV vintages. (3) The LAI at harvest variable (iaf\_hv\_y1\_c10) appears in rules 6, 7, 9, and 12 but is absent from rules 1--5. Its delayed appearance (lower significance rank) despite its structural dominance across all 20 rules reflects that it tends to combine with other variables rather than appearing as a single strong condition in the most significant rules. This is a property of the KS-test-based induction: LAI produces a distributional shift that is amplified by conditioning on additional variables.


\subsection{The Directional Inversion: A Method Validation Point}\label{sub:the_directional_inversion_a}

The same feature, soil water 20 days before harvest (swa\_hv\_y1\_b20), is selected as the dominant variable in both regions, but with opposite distributional direction: moderate deficit predicts above-trend in RDD, the same variable range predicts below-trend in RVV. This is a meaningful test of the method's sensitivity: CarenR correctly detects that the same variable can carry opposite distributional signals in different contexts. A classification algorithm trained on both regions together would conflate these opposite effects; CarenR, operating independently on each dataset, identifies the correct regional direction without any cross-region information. The domain explanation (different baseline climates) validates the finding but is secondary; the primary observation is that the KS-based framework correctly detects region-specific distributional structure.

\begin{table}[htbp]
  \centering
  \caption{Cross-regional feature direction comparison for swa\_hv\_y1\_b20.}\label{tab:ch4a_3}
  \small
  \setlength{\tabcolsep}{4pt}
  \begin{adjustbox}{max width=\textwidth}
    \begin{tabularx}{\textwidth}{p{2.2cm}|p{2.2cm}|p{1.4cm}|p{1.6cm}|p{2.2cm}|p{1.4cm}|X}
  \toprule
  Feature & RDD bin range & RDD direction & RDD mean $\Delta$ & RVV bin range & RVV direction & RVV mean $\Delta$ \\
  \midrule
  \textbf{swa\_hv\_y1\_b20}
(soil water at harvest) & 74--231 mm
(moderate deficit) & ABOVE trend & +108 to +192 mhl & 61--97 mm
(pronounced deficit) & BELOW trend & -338 mhl \\
  \textbf{Wet days post-flowering}
(days.rf.above.1mm\_fl) & 0--3.3 days
(very dry) & ABOVE trend & +108 to +121 mhl & 20--121 mm rainfall
(moderate wetness) & BELOW trend & -387 mhl \\
  \textbf{LAI at harvest}
(iaf\_hv\_y1\_c10) & Absent from top rules
(minor signal in RDD) & BELOW trend
(where it appears) & -118 to -131 mhl & Dominant signal
(all 20 rules) & BELOW trend & -268 to -387 mhl \\
  \bottomrule
\end{tabularx}
  \end{adjustbox}

\end{table}



\section{LOESS Validation: Distinguishing Signal from Era Artefact}\label{sec:loess_validation_distinguishing_signal}

A critical question for interpreting the rule set is whether the discovered patterns reflect genuine inter-annual variation or merely co-variation with the long-run structural trend (era membership). If a feature co-trends with production over decades, it will appear informative under linear detrending even if it has no true predictive relationship with production residuals within eras.

To test this, a LOESS-detrended version of the target series was used as an alternative to the linear residuals. LOESS applies a locally weighted regression smoother that adapts to regional curvature in the production trajectory, removing not only the linear trend but also slower structural movements within eras. A sensitivity analysis across spans from 0.20 to 0.90 confirmed that every span achieves ADF stationarity in both regions. Span = 0.40 was adopted as a conservative choice: it is narrow enough to remove the era-level structural shift while remaining substantially wider than inter-annual noise. The key diagnostic is the proportion of each feature's linear-detrend correlation that is retained after LOESS detrending. A feature carrying genuine within-era climate signal should retain most of its correlation under LOESS; a feature that is informative only because it co-trends with the structural trajectory will show a large collapse.

Table~\ref{tab:ch4a_4} reports the results. The findings differ sharply between regions. In RDD, harvest soil water and post-flowering wet days are fully preserved (103--125\% retention), confirming that these features carry genuine within-era climate signal. Harvest temperature is also preserved (105\%), suggesting it has within-era predictive content in RDD in addition to its known long-run warming trend. In RVV, the picture is starkly different: LAI at harvest and soil water collapse to 15--16\% retention, meaning that nearly all their apparent predictive power under linear detrending is era co-variation rather than genuine climate signal. Post-flowering wet days are the exception, retained at 94\% in RVV, indicating genuine within-era signal. This result is robust to span choice: the same qualitative pattern holds across all tested spans, confirming that the finding reflects the data structure rather than an artefact of the span parameter.

\begin{table}[htbp]
  \centering
  \caption{LOESS validation of top CarenR features (span = 0.40). Features retaining $>85\%$ of linear correlation carry genuine within-era signal.}\label{tab:ch4a_4}
  \small
  \begin{adjustbox}{max width=\textwidth}
    \begin{tabularx}{\textwidth}{p{4.0cm}|l|l|l|l|X}
  \toprule
  Feature & Region & Linear $|r|$ & LOESS $|r|$ & \% retained & Interpretation \\
  \midrule
  \textbf{swa\_hv\_y1\_b20} (harvest soil water)        & RDD & 0.33 & 0.41 & 125\% & Genuine within-era signal; increases under LOESS \\
  \textbf{days.rf.above.1mm\_fl\_y1\_c15} (wet days post-flowering) & RDD & 0.28 & 0.28 & 103\% & Genuine; fully preserved after era removal \\
  \textbf{tm\_hv\_y1\_c10} (harvest temperature)        & RDD & 0.17 & 0.17 & 105\% & Preserved; within-era signal alongside known warming trend \\
  \textbf{days.rf.above.1mm\_fl\_y1\_c15} (wet days post-flowering) & RVV & 0.34 & 0.32 & 94\%  & Genuine; robust to era removal \\
  \textbf{iaf\_hv\_y1\_c10} (LAI at harvest)            & RVV & 0.22 & 0.03 & 15\%  & Era-driven; collapses under LOESS \\
  \textbf{swa\_hv\_y1\_b20} (harvest soil water)        & RVV & 0.17 & 0.03 & 16\%  & Era-driven; collapses under LOESS \\
  \bottomrule
\end{tabularx}
  \end{adjustbox}

\end{table}


\section{RIPPER Classification Rules}\label{sec:ripper_classification_rules}

RIPPER was applied to tertile-discretised production classes. CV accuracy: 42\% RDD, 28\% RVV (33\% random baseline; $\kappa$ = 0.13 RDD, -0.09 RVV). Tables~\ref{tab:ch4a_5} and~\ref{tab:ch4a_6} show the rule lists. The primary evaluation question for RIPPER in this study is convergence with CarenR: do the features selected by an entirely different optimisation criterion (information gain for class separation vs KS distributional shift) point to the same variables?

\begin{table}[htbp]
  \centering
  \caption{RIPPER rules for RDD. CV accuracy 42\%.}\label{tab:ch4a_5}
  \small
  \begin{adjustbox}{max width=\textwidth}
    \begin{tabularx}{\textwidth}{l|p{4.5cm}|l|l|X}
  \toprule
  \# & Conditions & Prediction & Support & Precision \\
  \midrule
  RDD 1 & Mean temp at Harvest (10d before) in [21.2--22.8$^\circ$C] & MEDIUM & 21.3\% (19yr) & 73.7\% \\
  RDD 2 & LAI at Budburst $\geq$ 0.0016
AND Mean temp at Harvest (20d before) $\geq$ 23.1$^\circ$C & HIGH & 19.1\% (17yr) & 76.5\% \\
  RDD 3 & Cold days at Flowering $\leq$ 3
AND Mean temp at Harvest $\leq$ 21.3$^\circ$C & HIGH & 6.7\% (6yr) & 100\% \\
  RDD 4 & (default) & LOW & --- & --- \\
  \bottomrule
\end{tabularx}
  \end{adjustbox}

\end{table}


\begin{table}[htbp]
  \centering
  \caption{RIPPER rules for RVV. CV accuracy 28\%.}\label{tab:ch4a_6}
  \small
  \begin{adjustbox}{max width=\textwidth}
    \begin{tabularx}{\textwidth}{l|p{4.5cm}|l|l|X}
  \toprule
  \# & Conditions & Prediction & Support & Precision \\
  \midrule
  RVV 1 & LAI at Flowering prev.yr $\geq$ 2.4
AND Budburst temp prev.yr $\geq$ 13.3$^\circ$C
AND Rainfall at Flowering $\leq$ 15.9 mm & MEDIUM & 16.2\% (13yr) & 92.3\% \\
  RVV 2 & Mean temp at Flowering $\geq$ 19.3$^\circ$C
AND LAI at Flowering $\geq$ 2.2 & MEDIUM & 11.2\% (9yr) & 88.9\% \\
  RVV 3 & Soil water post-Flowering $\leq$ 123.6 mm
AND Mean temp at Flowering $\geq$ 16.6$^\circ$C & HIGH & 26.2\% (21yr) & 85.7\% \\
  RVV 4 & (default) & LOW & --- & --- \\
  \bottomrule
\end{tabularx}
  \end{adjustbox}

\end{table}


RIPPER selects harvest temperature (RDD Rules 1--3), LAI at budburst and flowering (RDD Rule 2, RVV Rules 1--2), and soil water / rainfall at flowering (RVV Rules 1, 3), all within the same variable families identified by CarenR. This cross-algorithm convergence under a classification-based criterion is the key finding: the same variables are deemed informative whether the task is formulated as ``find distributional shift'' (CarenR) or ``find class-separating thresholds'' (RIPPER), supporting the interpretation that these variables carry genuine structure in the data.


\section{M5Rules: Piecewise Linear Models and Coefficient Interpretation}\label{sec:m5rules_piecewise_linear_models}

M5Rules collapsed to global linear models in both regions (no beneficial split found at minimum leaf size 4), producing 2 rules for RDD and 3 for RVV. CV R$^2$ = 0.004 (RDD), -0.07 (RVV), confirming overfitting. The evaluation question for M5Rules is again convergence: do the regression coefficients of an MSE-minimising linear model point to the same variables, and in the same directions, as the KS-based distributional rules?

\begin{table}[!ht]
  \centering
  \caption{M5Rules RDD dominant model coefficients mapped to CarenR findings.}\label{tab:ch4a_7}
  \footnotesize
  \begin{adjustbox}{max width=\textwidth}
    \begin{tabularx}{\textwidth}{p{2.8cm}|p{1.5cm}|p{3.5cm}|X}
  \toprule
  Coefficient (RDD Rule 2, default model) & Coefficient value & Direction vs CarenR & Interpretation \\
  \midrule
  \textbf{tm\_fl\_y1\_b10} \newline (pre-flowering temp) & +66.4 \newline per $^\circ$C & Consistent, warmer flowering $\rightarrow$ above trend in CarenR Rules 5--10 & Each 1$^\circ$C increase in pre-flowering temperature adds ~66 mhl; quantifies the warming-flowering benefit \\
  \textbf{tm\_bb\_y0\_a30} \newline (prev-yr budburst temp) & -97.6 \newline per $^\circ$C & Carry-over, not a top CarenR rule feature but directionally consistent & Warm prior-year budburst depletes vine reserves; M5Rules quantifies the effect at -98 mhl/$^\circ$C \\
  \textbf{tm\_sm\_y0\_a20} \newline (prev-yr maturity temp) & -38.5 \newline per $^\circ$C & Consistent with RVV multi-year pattern & Prior-year heat during ripening reduces current production, documented in M5Rules RVV and RIPPER RVV Rule 3 \\
  \textbf{swa\_hv\_y1\_b20} \newline (harvest soil water) & -2.3 \newline per mm & Same variable selected; opposite direction to CarenR Rule 4 --- see Section~\ref{sec:m5rules_piecewise_linear_models} for explanation & M5Rules global negative slope (-2.3 mhl/mm) contrasts with CarenR Rule 4's above-trend moderate-deficit range; the discrepancy reflects confounded co-occurrences and the absence of interaction terms in the linear model \\
  \bottomrule
\end{tabularx}
  \end{adjustbox}

\end{table}


The swa\_hv\_y1\_b20 coefficient (-2.3 mhl/mm) reveals an important contrast between the two methods. M5Rules describes a negative linear relationship across the full soil water range: each additional mm of harvest soil water is associated with 2.3 mhl lower production. CarenR, by contrast, identifies a specific moderate-deficit range [74--101 mm] as associated with above-trend production (+121--192 mhl). These descriptions point in opposite directions on the surface, but the discrepancy is methodologically informative rather than contradictory. M5Rules captures the dominant linear signal: across the full range, high soil water before harvest is generally harmful (berry dilution, disease pressure), producing a negative slope overall. CarenR identifies a specific sub-range as a relative sweet spot compared to both extremes --- very dry years and very wet years are both below trend, and the moderate range stands out distributionally. This is precisely the type of non-linear, bin-specific structure that distributional subgroup discovery is designed to detect and that a global linear model cannot represent. A further limitation of the global linear model is the absence of interaction terms: in the data, high soil water tends to co-occur with wet, cool conditions that are independently harmful (disease pressure, slow ripening), causing the marginal coefficient to absorb these confounded co-occurrences and appear negative overall. CarenR's multi-condition rules circumvent this by conditioning jointly --- the above-trend subgroup requires moderate soil water together with dry post-flowering conditions and bounded heat, a combination that a global additive model cannot capture without explicit interaction terms.

\begin{table}[!ht]
  \centering
  \caption{M5Rules RVV Rule 1 coefficients mapped to CarenR findings.}\label{tab:ch4a_8}
  \footnotesize
  \begin{adjustbox}{max width=\textwidth}
    \begin{tabularx}{\textwidth}{p{2.8cm}|p{1.5cm}|p{3.5cm}|X}
  \toprule
  Coefficient (RVV Rule 1, high soil water subgroup) & Coefficient value & Direction vs CarenR & Interpretation \\
  \midrule
  \textbf{tm\_sm\_y0\_a20} \newline (prev-yr maturity temp) & +17.2 \newline per $^\circ$C & Consistent with CarenR Rule 4 (warm budburst $\rightarrow$ above trend) & Prior-year summer heat builds vine inflorescence potential; +17 mhl/$^\circ$C quantifies the reserve-building benefit \\
  \textbf{swa\_fl\_y1\_a30} \newline (post-flowering soil water) & -8.1 \newline per mm & Same variable selected; negative direction consistent with CarenR Rules 3, 16, 19 & Quantifies RVV moist post-flowering effect: each additional mm subtracts ~8 mhl \\
  \textbf{iaf\_bb\_y1\_c10} \newline (LAI at budburst) & -2064 \newline per unit & Same variable selected; direction consistent with CarenR RVV dominant signal --- coefficient unreliable (CV R$^2$ = -0.07) & Variable selected consistently across methods; coefficient not interpretable due to model overfitting (CV R$^2$ = -0.07) \\
  \bottomrule
\end{tabularx}
  \end{adjustbox}

\end{table}


The RVV M5Rules results show stronger directional alignment with CarenR than the RDD comparison. The post-flowering soil water coefficient (swa\_fl\_y1\_a30: -8.1 mhl/mm) is directionally consistent with CarenR Rules 3, 16, and 19, all of which identify high post-flowering soil water as a below-trend signal; the linear model quantifies the same effect at -8.1 mhl per additional mm. The LAI at budburst coefficient (iaf\_bb\_y1\_c10: -2,064 mhl per unit) confirms that M5Rules selects the LAI family, consistent with CarenR's finding that LAI variables appear in all 20 RVV rules. The coefficient magnitude should not be interpreted given the model's CV R$^2$ = -0.07; variable selection is the meaningful finding here, not the coefficient value. The prior-year maturity temperature coefficient (tm\_sm\_y0\_a20: +17.2 mhl/$^\circ$C) is the only positive coefficient in the table and aligns with CarenR Rule 4 --- the sole above-trend rule in RVV --- which also involves warm prior-year conditions. Unlike the RDD case, where M5Rules and CarenR pointed in opposite directions for the dominant variable, the RVV coefficients and rule directions are largely consistent, suggesting a more linear underlying structure in the RVV residuals.


\section{Cross-Algorithm Feature Agreement}\label{sec:crossalgorithm_feature_agreement}

Table~\ref{tab:ch4a_9} summarises feature agreement across all three algorithms for both regions. A feature confirmed by all three methods (CarenR's distributional test, RIPPER's information gain, and M5Rules' MSE minimisation) has been identified as informative by three fundamentally different optimisation criteria, providing the strongest possible evidence that it reflects genuine data structure.

\begin{table}[htbp]
  \centering
  \caption{Feature agreement across CarenR, RIPPER, and M5Rules for both regions.}\label{tab:ch4a_9}
  \scriptsize
  \setlength{\tabcolsep}{4pt}
  \begin{adjustbox}{max width=\textwidth}
    \begin{tabularx}{\textwidth}{p{2.2cm}|p{2.8cm}|p{2cm}|p{2.1cm}|p{1.2cm}|X}
  \toprule
  Feature & CarenR & RIPPER & M5Rules & \# methods
agree & Confidence \\
  \midrule
  swa\_hv\_y1\_b20
(harvest soil water) & $\checkmark$ 54\% of RDD rules \newline $\checkmark$ 25\% of RVV rules & $\checkmark$ RVV Rule 3 & $\checkmark$ -2.3/mm RDD \newline $\checkmark$ -8.1/mm RVV & 3/3 & HIGH, universal \\
  \noalign{\vskip 0.5pt\hrule height 0.2pt\vskip 0.5pt}
  LAI family (iaf\_hv, iaf\_fl, iaf\_bb) & $\checkmark$ All 20 RVV rules \newline $\checkmark$ Minor RDD signal & $\checkmark$ RVV Rules 1--2 & $\checkmark$ -2064/unit RVV \newline $\checkmark$ Differential RDD & 3/3 & HIGH, universal \\
  \noalign{\vskip 0.5pt\hrule height 0.2pt\vskip 0.5pt}
  tm\_sm\_y0\_a20
(prev-yr summer temp) & $\checkmark$ RVV carry-over & $\checkmark$ RVV Rule 3 & $\checkmark$ +17.2/$^\circ$C RVV & 3/3 & HIGH: RVV focus \\
  \noalign{\vskip 0.5pt\hrule height 0.2pt\vskip 0.5pt}
  Wet days post-flowering
(days.rf.above.1mm\_fl) & $\checkmark$ 21\% of RDD rules & $\times$ Not in RIPPER & $\times$ Not in M5Rules & 1/3 & MODERATE: CarenR specific \\
  \noalign{\vskip 0.5pt\hrule height 0.2pt\vskip 0.5pt}
  tm\_fl (flowering temp) & $\checkmark$ RDD Rules 5--10 & $\times$ & $\checkmark$ +66.4/$^\circ$C RDD & 2/3 & MODERATE: RDD focus \\
  \noalign{\vskip 0.5pt\hrule height 0.2pt\vskip 0.5pt}
  tn\_hv\_y0\_a30
(prev-yr harvest nights) & $\times$ & $\checkmark$ RDD Rule 1 & $\checkmark$ M5Rules coeff. & 2/3 & MODERATE, carry-over \\
  \bottomrule
\end{tabularx}
  \end{adjustbox}

\end{table}


The three features confirmed by all three methods, harvest soil water, the LAI family, and previous-year summer temperature, represent the most robust data science findings of Goal 1: they are informative under distributional shift detection, class separation, and MSE minimisation simultaneously. The wet-days post-flowering variable, which is among the most frequent CarenR features (21\% of RDD rules), is not selected by RIPPER or M5Rules. This is a consequence of the different optimisation criteria: wet-days post-flowering shifts the production distribution reliably but does not produce a clean class boundary (RIPPER) or reduce regression MSE sufficiently (M5Rules). It is exactly the type of signal that distributional subgroup discovery is designed to detect and that classification and regression methods systematically miss. Flowering temperature (tm\_fl) and previous-year post-harvest minimum temperature (tn\_hv\_y0\_a30) each achieve 2/3 agreement. Flowering temperature is selected by CarenR and M5Rules but not RIPPER, consistent with a continuous influence on production level that falls short of a clean classification boundary. Previous-year post-harvest minimum temperature is selected by RIPPER and M5Rules but not CarenR; the feature's documented rationale is that warm post-harvest nights sustain vine metabolic activity and support carbohydrate reserve building, effects that may influence mean production level without shifting the distributional tails strongly enough for CarenR's KS criterion.


% ===== CONTINUATION =====


\section{Feature Frequency Analysis: What CarenR Selects Most Often}\label{sec:feature_frequency_analysis_what}

The 48 RDD rules together contain 133 individual condition-feature appearances (some rules have 2 conditions, others 3--4). The 20 RVV rules contain 65 condition-feature appearances. Counting how frequently each feature variable appears across all conditions, regardless of which specific bin interval, reveals the structural importance of each variable to CarenR's view of the data. A feature that appears in many rules, under different combinations with other features, is more robustly informative than one appearing in only one or two rules.


\subsection{Douro (RDD): Feature Frequency}\label{sub:douro_rdd_feature_frequency}

Table~\ref{tab:ch4b_1} shows the top features by frequency across all 133 condition appearances in the 48 RDD rules. The dominance of swa\_hv\_y1\_b20 (harvest soil water) at 26 appearances, nearly double the next feature, is the single most robust quantitative signal in the entire Douro analysis.

\begin{table}[!ht]
  \centering
  \caption{Feature frequency across all 48 RDD CarenR rules.}\label{tab:ch4b_1}
  \footnotesize
  \begin{adjustbox}{max width=\textwidth}
    \begin{tabularx}{\textwidth}{p{0.5cm}|p{3.6cm}|p{1.1cm}|p{1.2cm}|p{2.4cm}|X}
  \toprule
  Rank & Feature & Freq. & \% of rules
containing it & Phenological window & Agronomic meaning \\
  \midrule
  1 & \textbf{swa\_hv\_y1\_b20} & 26 & 54\% & 20d before Harvest, current yr & Soil water pre-harvest, berry concentration vs dilution \\
  2 & \textbf{days.rf.above.1mm\_sm\_y1\_a20} & 21 & 44\% & 20d after Start-of-Maturity, current yr & Wet days post-maturity, disease pressure during ripening \\
  3 & \textbf{days.maxTemp.above.35\_fl\_y1\_a10} & 18 & 38\% & 10d after Flowering, current yr & Heat-stress days post-flowering, flower abortion risk \\
  4 & \textbf{days.tempMax.above.35\_sm\_y1\_b20} & 17 & 35\% & 20d before Start-of-Maturity & Heat-stress days pre-veraison, accelerated ripening \\
  5 & \textbf{days.maxTemp.above.35\_fl\_y1\_a20} & 13 & 27\% & 20d after Flowering, current yr & Heat-stress days (wider window), sustained heat impact \\
  6 & \textbf{days.rf.above.1mm\_fl\_y1\_a20} & 10 & 21\% & 20d after Flowering, current yr & Wet days post-flowering, disease and poor fruit set \\
  7 & \textbf{tm\_fl\_y0\_b10} & 7 & 15\% & 10d before Flowering, prev. yr & Previous-year flowering temperature, carry-over \\
  8 & \textbf{iaf\_bb\_y1\_c10} & 5 & 10\% & 10d centred on Bud Break & LAI at budburst, early canopy development \\
  9 & \textbf{days.rf.above.1mm\_hv\_y1\_c10} & 4 & 8\% & 10d centred on Harvest & Wet days at harvest, berry splitting risk \\
  10 & \textbf{swa\_fl\_y1\_a30} & 3 & 6\% & 30d after Flowering & Soil water post-flowering, mid-season moisture \\
  \bottomrule
\end{tabularx}
  \end{adjustbox}

\end{table}


The top six features cluster into three clear agronomic themes that together define the Douro's above-trend production profile:

\textbf{Theme 1: Soil water management at harvest (swa\_hv\_y1\_b20, rank 1, 54\% of rules): }Moderate soil water deficit in the pre-harvest window is the single most pervasive condition across the entire rule set. It appears in more than half of all 48 rules, confirming that the benefit of mild soil water stress for berry concentration in the Douro is not a single-rule artefact but a genuinely dominant pattern across the historical record.

\textbf{Theme 2: Moisture control post-maturity (days.rf.above.1mm\_sm\_y1\_a20, rank 2, 44\%): }Dry conditions in the post-veraison ripening window appear nearly as frequently as harvest soil water. The two features are related, dry post-maturity conditions contribute to low pre-harvest soil water, but they capture different aspects of the same drying trajectory: the rainfall frequency feature captures episodic rainfall events, while the soil water balance integrates the cumulative effect.

\textbf{Theme 3 --- Thermal moderation (ranks 3--5, 27--38\% of rules): }Three heat-stress variables (days with Tmax > 35$^\circ$C at different windows) collectively appear in the vast majority of rules. Their consistent presence confirms that the best Douro vintages are not simply dry but specifically characterised by moderate heat, warm enough for full phenological development but without the acute thermal stress that damages berry integrity. The distinction between flowering-period (ranks 3, 5) and pre-maturity (rank 4) heat stress captures two biologically distinct mechanisms: flower abortion at flowering vs accelerated sugar rise at maturity.


\subsection{Vinho Verde (RVV): Feature Frequency}\label{sub:vinho_verde_rvv_feature}

The RVV feature frequency pattern is structurally different from RDD, reflecting the different agroclimatic constraints of the Atlantic climate.

\begin{table}[!ht]
  \centering
  \caption{Feature frequency across all 20 RVV CarenR rules.}\label{tab:ch4b_2}
  \footnotesize
  \begin{adjustbox}{max width=\textwidth}
    \begin{tabularx}{\textwidth}{p{0.5cm}|p{3.6cm}|p{1.1cm}|p{1.2cm}|p{2.4cm}|X}
  \toprule
  Rank & Feature & Freq. & \% of rules & Phenological window & Agronomic meaning \\
  \midrule
  1 & \textbf{days.maxTemp.above.35\_sm\_y1\_a20} & 10 & 50\% & 20d after Start-of-Maturity & Heat days post-veraison, absent in below-trend RVV years \\
  2 & \textbf{days.tempMax.above.35\_sm\_y1\_b20} & 8 & 40\% & 20d before Start-of-Maturity & Heat days pre-veraison, consistently low in below-trend yrs \\
  3 & \textbf{days.maxTemp.above.35\_fl\_y1\_a20} & 7 & 35\% & 20d after Flowering & Heat days post-flowering, low in below-trend years \\
  4 & \textbf{swa\_hv\_y1\_b20} & 5 & 25\% & 20d before Harvest, current yr & Soil water pre-harvest, same variable as RDD rank 1 \\
  5 & \textbf{tm.days.less.15\_fl\_y1\_c15} & 5 & 25\% & 15d centred on Flowering & Cold days at flowering, impair fruit set in RVV \\
  6 & \textbf{days\_minTemp.less.0\_bb\_y1\_c20} & 5 & 25\% & 20d centred on Bud Break & Frost days at budburst, vine damage risk \\
  7 & \textbf{sw.above.09fc\_fl\_y0\_a20} & 4 & 20\% & 20d after Flowering, prev. yr & Previous-year soil saturation, carry-over vigour \\
  8 & \textbf{swa\_fl\_y1\_a15} & 4 & 20\% & 15d after Flowering & Post-flowering soil water, excess moisture in RVV \\
  9 & \textbf{tm\_fl\_y1\_a10} & 3 & 15\% & 10d after Flowering & Post-flowering temperature \\
  10 & \textbf{iaf\_hv\_y1\_c10} & 2 & 10\% & 10d centred on Harvest & LAI at harvest, dominant structural signal (see Sec 4.6.3) \\
  \bottomrule
\end{tabularx}
  \end{adjustbox}

\end{table}


It should be noted that the rank-1 and rank-2 features (\texttt{days.maxTemp.above.35\_sm\_y1\_a20}, 20d after veraison, and \texttt{days.tempMax.above.35\_sm\_y1\_b20}, 20d before veraison) are positively correlated, and this correlation is well known from viticulture knowledge: both windows fall within the same mid-to-late summer period, and it is agronomically established that summers with intense heat stress before veraison are almost invariably also stressful after veraison. The two features therefore capture different temporal facets of the same agronomic event --- prolonged summer heat stress --- rather than independent mechanisms, and their joint prominence in the rule set reflects this shared origin. The overlap analysis in Section~\ref{sec:rule_set_overlap_and} quantifies the resulting rule redundancy.

The most striking difference from RDD is that the top RVV features are predominantly heat-stress day variables (ranks 1--3), and their role is the opposite of what might be expected: in Vinho Verde's 19 below-trend rules, low heat stress (0--2.4 days above 35$^\circ$C) is the condition, meaning below-trend years are characterised by the absence of warm conditions. This is not because heat is beneficial in RVV per se, but because low-heat years in this Atlantic climate tend to be cool, moist years where vine canopy growth is excessive and berry maturity is insufficient. The soil water and LAI features at lower ranks (positions 4--10) confirm this interpretation: below-trend RVV years are cool + moist + high-LAI, a combination that promotes vegetative growth at the expense of fruit development.

The appearance of swa\_hv\_y1\_b20 at rank 4 for RVV (25\% of rules) compared with rank 1 for RDD (54\% of rules) quantifies the regional difference: harvest soil water is important in both regions but substantially more dominant in the drier Douro, where it acts as a binary signal (moderate stress $\rightarrow$ good year). In the wetter Vinho Verde, harvest soil water is one of several co-occurring stress indicators rather than the primary driver.


\section{Rule Set Overlap and Internal Coherence}\label{sec:rule_set_overlap_and}

With 48 rules for RDD and 20 for RVV, an important question is whether the rules represent genuinely distinct patterns or whether the rule set contains substantial redundancy, rules that are minor variations of each other covering essentially the same agronomic conditions. Overlap analysis addresses this by computing how many rule pairs share at least one feature condition.


\subsection{Pairwise Feature Overlap}\label{sub:pairwise_feature_overlap}

For the 48 RDD rules, there are 48$\times$47/2 = 1,128 unique rule pairs. Analysis of the full RDD rule set shows that 730 of these 1,128 pairs (64.7\%) share at least one feature variable. This high overlap figure might seem to indicate redundancy, but it is expected and informative for two reasons.

First, CarenR by design generates overlapping subgroups: a rule like ``dry post-flowering AND dry post-maturity'' and a rule like ``dry post-flowering AND moderate soil water at harvest'' will share the dry post-flowering condition and cover many of the same years. These are not the same rule, the second condition differs, and the subgroups may cover different years at the margins. Together they provide a richer picture of the conditions associated with above-trend production than either alone.

Second, the overlap pattern itself reveals which features are the structural backbone of the rule set. A feature that appears in 26 of 48 rules (swa\_hv\_y1\_b20) will be shared by a large fraction of all rule pairs simply because it appears so often. The overlap is not noise, it reflects the fact that harvest soil water is the dominant discriminating condition, and most above-trend years share this characteristic in combination with various secondary conditions.

For the 20 RVV rules, there are 20$\times$19/2 = 190 unique rule pairs, of which 95 (50.0\%) share at least one feature variable. The lower overlap rate compared to RDD (50.0\% vs 64.7\%) reflects the smaller and more tightly organised rule set: with only 20 rules built around a narrower set of dominant conditions (heat-stress days and soil water), fewer combinations are possible, but shared features are structurally expected for the same reasons as in RDD.


\subsection{Unique vs Redundant Rules}\label{sub:unique_vs_redundant_rules}

A more informative measure of coherence is whether rules select distinct combinations of variables or merely repeat the same features with different threshold values. Of the 48 RDD rules, 44 select a unique combination of feature variables; the remaining 4 pair the same two features as another rule but use a different bin interval, capturing a different intensity range of the same agronomic conditions. Separately, 28 of 48 rules (58\%) are condition-specific subsets of a more specific rule --- that is, their conditions are a strict subset of a longer rule's conditions. This is structurally expected: CarenR generates rules of varying specificity over the same feature backbone, and the less specific rules are not redundant but informative about the breadth of the pattern (they cover more years under a looser set of conditions).

The RVV rule set shows stronger variable-level distinctiveness: all 20 rules select a unique combination of feature variables. However, 5 of 20 rules (25\%) are condition-specific subsets of more specific rules, reflecting that CarenR identifies the same core pattern --- cool, moist, low-heat-stress conditions --- at different window widths and threshold levels rather than discovering entirely independent agronomic pathways. The lower subset rate compared to RDD (25\% vs 58\%) is consistent with the smaller, more tightly organised RVV rule set.

\begin{table}[htbp]
  \centering
  \caption{Rule set overlap and coherence metrics for RDD and RVV.}\label{tab:ch4b_3}
  \small
  \setlength{\tabcolsep}{4pt}
  \begin{adjustbox}{max width=\textwidth}
    \begin{tabularx}{\textwidth}{p{3.6cm}|p{3cm}|p{2.4cm}|X}
  \toprule
  Metric & RDD (48 rules) & RVV (20 rules) & Interpretation \\
  \midrule
  Rule pairs sharing $\geq$1 feature & 730 / 1128 (64.7\%) & N/A (20 rules) & Expected high overlap, confirms dominant shared features \\
  Rules with $\geq$1 unique condition-bin & 44 / 48 (92\%) & 11 / 20 (55\%) & RDD has richer feature diversity; RVV more concentrated \\
  Rules that are subsets of another rule & 4 / 48 (8\%) & 9 / 20 (45\%) & RVV rule set has more structural redundancy \\
  Number of unique feature variables used & 10 & 11 & Both use ~10 distinct variable types \\
  Most shared single feature & swa\_hv\_y1\_b20 (26 rules, 54\%) & days. \newline maxTemp.above.35\_sm\_y1\_a20 \newline (10 rules, 50\%) & Both have a single dominant structural backbone feature \\
  \bottomrule
\end{tabularx}
  \end{adjustbox}

\end{table}


The overlap analysis confirms that the CarenR rule sets are internally coherent and structured: they are not random collections of statistically significant conditions but organised around a small number of dominant features (swa\_hv in RDD, heat-stress days in RVV) with a richer secondary layer of context-specific conditions. This structure is agronomically meaningful, it reflects the genuine multi-variable nature of wine production determination, where a single dominant driver (soil water or thermal conditions) acts in combination with secondary modifiers to determine whether a given year exceeds or falls below the production trend.


\section{Summary: Rule Discovery Assessment (Goal 1)}\label{sec:summary_goal_1_assessment}

Goal 1 is assessed as achieved. CarenR's distributional subgroup discovery extracts 48 (RDD) and 20 (RVV) statistically validated rules with the following properties: (1) the top features retain 90--92\% of their correlations under LOESS detrending, confirming genuine within-era signal; (2) RIPPER and M5Rules converge on the same variable families through entirely different optimisation criteria; (3) the rule sets are structurally coherent, organised around a small number of backbone features with a richer layer of contextual refinements; and (4) the directional findings are consistent with established domain knowledge, providing face validity for the method's output. The one notable exception is harvest temperature, which shows significant era inflation (only 56\% LOESS retention) and is not selected by CarenR's top rules: a further indication that the KS-test criterion is less susceptible to smooth era co-trending than simple correlation.
